\subsection{Two elementary lemmas about low-degree polynomials}\label{ht:two-elementary-lemmas-about-low-degree-polynomials}

For \(t \in [-1,1]\), let \(P_t\) be the product spin distribution with mean \(\mathbb E X_i=t\). If \(h\) is a degree-\(k\) multilinear polynomial, then

\[
q_h(t):=\mathbb E_{P_t}h
=\sum_{S\subseteq[n]}\widehat h(S)t^{|S|}
\]

is a univariate polynomial of degree at most \(k\). If \(|h|\le 1\) on the cube, then \(|q_h(t)|\le 1\) throughout \([-1,1]\).

\textbf{Lemma HT3 (symmetric mixture comparison).} Suppose a measure \(\mu\) is a mixture \(\mu=\mathbb E[P_T]\), with a symmetric random variable \(T \in [-1,1]\). For cube-bounded \(h\) of degree at most \(k\),

\[
|\mathbb E_\mu h-\mathbb E_Uh|
\le \tfrac12 k^4\,\mathbb E T^2.
\tag{HT3}
\]

For \(k=0\), the right side and left side both vanish.

\textbf{Proof.} Markov's polynomial inequality, applied twice, gives \(\lVert q_h''\rVert_\infty\le k^4\). Taylor's theorem gives

\(|q_h(t)-q_h(0)-t q_h'(0)| \le  k^4 t^2/2\).

The linear term integrates to zero because \(\mathbb E T=0\). \(\square\)

For \(f\) as above, applying HT3 to \(f\) and the multilinear representation of \(f^2\) gives

\[
|\operatorname{Var}_\mu f-\operatorname{Var}_U f|
\le 9d^4\,\mathbb E T^2.
\tag{4.1}
\]

Indeed, the second-moment discrepancy is at most \(8 d^4 \mathbb E T^2\), and the squared-mean discrepancy is at most \(d^4 \mathbb E T^2\) because both means lie in \([-1,1]\). Multilinearization of \(f^2\) does not increase its degree above \(2d\) and preserves its values on the cube.

\textbf{Lemma HT4 (central density preserves influence).} Suppose the mixing variable \(T\) has density at least \(\eta/a\) on \([-a,a]\), where \(0<a\le 1\) and \(\eta>0\). Then

\[
I_i^\mu(f)\ge \frac{\eta}{16d^2}I_i^U(f)
\quad\text{for every }i.
\tag{HT4}
\]

\textbf{Proof.} Put \(h(t)=\mathbb E_{P_t}(D_i f)^2\). This is a nonnegative polynomial on \([-1,1]\) of degree at most \(k=2(d-1)\). Set \(K=\max(1,k)\le 2d\) and \(M=\max_{[-a,a]}h\). The rescaled Markov inequality gives

\(\max_{[-a,a]}|h'| \le  K^2 M/a\).

At a maximizer of \(h\), there is a one-sided interval inside \([-a,a]\) of length at least \(a/(2K^2)\) on which \(h\ge M/2\). Therefore

\[
\int_{-a}^a h(t)dt\ge\frac{aM}{4K^2}
\ge\frac{a h(0)}{16d^2}.
\]

Integrating against the density lower bound proves HT4. This proof also covers constant \(h\) via \(K=1\). \(\square\)

The classical Markov inequality used here is \(\lVert q'\rVert_{[-1,1]}\le k^2\lVert q\rVert_{[-1,1]}\). One primary mathematical reference stating it is \href{https://wayback.cecm.sfu.ca/~pborwein/PAPERS/P27.pdf}{Borwein, \emph{Markov's inequality for polynomials with real zeros}, Proc. AMS 93 (1985), equation (1)}.
